\documentclass{article}
\usepackage{amsmath,amssymb}
\usepackage[margin=1in]{geometry}
\allowdisplaybreaks
\usepackage[UTF8]{ctex}
\begin{document}

\section*{仿射近似动态规划：(D1)}

\begin{equation}
(\mathbf{D1})\qquad
\min_{\theta,V}\;\theta_1+\sum_i V_{1,i}c_i
\tag{6}
\end{equation}
受…的约束
\begin{equation}
\begin{aligned}
&\theta_t-\theta_{t+1}
 +\sum_i\left(V_{t,i}x_i-V_{t+1,i}\left(x_i-\sum_{j\in S}\lambda P_j(S)a_{ij}\right)\right)
 \geq \sum_{j\in S}\lambda P_j(S)f_j,\\
&\hspace{25mm}\forall\,t=1,\ldots,\tau,\quad x\in X,\quad S\subseteq N(x).
\end{aligned}
\tag{7}
\end{equation}
末端系数已经确定：
\[
\theta_{\tau+1}=0,\qquad V_{\tau+1,i}=0\quad\forall\,i=1,\ldots,m.
\]
决策变量是$\theta_t$和$V_{t,i}$，它们分别用于$t=1,\ldots,\tau$和$i=1,\ldots,m$的运算。在打印的(D1)块中，这些决策变量没有受到任何符号限制。状态变量$x$以及提供集合$S$用于表示不等式；它们并不是这个线性程序的额外决策变量。

\section*{固定的输入和索引化的可行集合}

有$m$个航班班次和$n$种产品。产品指的是航班行程与票价等级的组合。初始容量向量和产品集合是……
\[
c=(c_1,\ldots,c_m),\qquad N=\{1,\ldots,n\},
\]
其中，$c_i$表示腿$i$的容量。产品$j$的运费为$f_j$。消耗矩阵就是$m\times n$矩阵。
\[
A\equiv(a_{ij}).
\]
其条目$a_{ij}$表示类$j$客户所需的资源$i$的整数数量。其第$i$行$A_i$是线路$i$的相关向量，而其第$j$列$A^j$则是产品$j$的相关向量。在整个过程中，$i$表示线路，$j$表示产品，而$t$表示时间；在(5)--(7)中，对$i$的所有求和运算都是针对$i=1,\ldots,m$进行的。

有$\tau$个离散的时段，按顺序排列。最后一个时段为$\tau$。在每个时段中，有$\lambda$的概率有一个客户到来，而概率为$1-\lambda$时则没有客户到来。在有一个客户到来并且有$S\subseteq N$种优惠方案的情况下，客户有$P_j(S)$的概率选择产品$j\in S$，并且没有购买行为。
\[
P_0(S)=1-\sum_{j\in S}P_j(S),\qquad S\subseteq N.
\]
选择概率是一个固定的输入值；在这个公式中并没有强制使用特定的选择模型。集合$S=\varnothing$表示没有提供任何产品，它属于可能的产品提供集合之一。

在一个周期的开始处，$x$表示未售出的座位数量，即$m$向量。状态空间与可行产品集合为：
\[
X=\{0,\ldots,c_1\}\times\cdots\times\{0,\ldots,c_m\},
\qquad
N(x)=\{j\in N:x\geq A^j\},\quad x\in X.
\]
因此，$N(x)$由可以在状态$x$下提供的产品组成；向量不等式要求每条路径都有足够的容量。初始状态为$c$。每个子集$S\subseteq N(x)$，包括空集在内，都对应着方程(7)中的一个约束条件。

\section*{仿射值函数构造}

设 $v_t(x)$ 为周期 $t,\ldots,\tau$ 结束时的最大总预期收入，此时周期 $t$ 开始时的状态为 $x$。 所选的近似值为
\begin{equation}
v_t(x)\approx\theta_t+\sum_i V_{t,i}x_i,
\qquad t=1,\ldots,\tau,\quad x\in X.
\tag{5}
\end{equation}
在这里，$V_{t,i}$估算了在时期$t$乘坐航班$i$的座位边际价值，而$\theta_t$是一个固定的偏移量。

继承下来的价值函数定义就是贝尔曼方程。
\begin{equation}
\begin{aligned}
v_t(x)
&=\max_{S\subseteq N(x)}\left\{
\sum_{j\in S}\lambda P_j(S)\bigl(f_j+v_{t+1}(x-A^j)\bigr)
 +(\lambda P_0(S)+1-\lambda)v_{t+1}(x)
\right\}\\
&=\max_{S\subseteq N(x)}\left\{
\sum_{j\in S}\lambda P_j(S)
\bigl[f_j-\bigl(v_{t+1}(x)-v_{t+1}(x-A^j)\bigr)\bigr]
\right\}+v_{t+1}(x),\\
&\hspace{30mm}\forall\,t=1,\ldots,\tau,\quad x\in X,
\end{aligned}
\tag{1}
\end{equation}
具有边界条件
\[
v_{\tau+1}(x)=0\qquad\forall\,x\in X.
\]
方程(1)定义了需要近似的精确值函数，而不是对仿射系数的额外约束。该论文通过将(5)代入(D0)中的贝尔曼不等式中，得到了(D1)；(7)则是由此得出的完整不等式体系，而上述固定的终端系数则实现了终端值的计算。

\paragraph{范围与来源。}
这是Zhang和Adelman的原文(D1)，标题为\emph{一种针对具有客户选择的网络收入管理的近似动态规划方法}，发表于《Transportation Science》43(3)，2009年，页码为~381--394，章节为~2.1（物理PDF第~4页），定义来自~1.1章节（物理PDF第~2--3页）。对偶问题(P1)未被选中。物理PDF第~9页的实施方案证实了定理~1所提出的单调性约束在数值实验中得到了应用；这些实验性补充并不包含在选定的印刷版公式(6)--(7)中。用于证明定理~1的在线附录并未提供，因此尚未验证；该附录中的任何条件均未在此处引用。

\end{document}
